\documentclass[../main.tex]{subfiles}
\begin{document}
\chapter{Introduction}
\section{Classical Mechanics}
\textit{Classical Mechanics} is based on two distinct concepts:
\begin{enumerate}
  \item \textbf{Particles} are ``point like'' objects with position given by $\vec{x}(t)$ and velocity $\dot{\vec{x}}(t)$.
  If we have an initial stating position and velocity $\vec{x}(t)$ and $\dot{\vec{x}}(t)$, then their motions are fully determined by Newton's Second Law:
  \[
    m \ddot{\vec{x}}(t) = F(\vec{x}(t), \dot{\vec{x}}(t))
  \]
  \item \textbf{Waves} are ``spread out'' objects described by $f(x, t)$ which is periodic in both $x$ and $t$.
  Their propagation is governed by the PDE:
  \[
    \pderiv[2]{f}{t} - c^2 \pderiv[2]{f}{x} = 0
  \]
  which, as we saw in IA Differential Equations, has solutions of the form:
  \[
    f_{\pm}(x, t) = A_{\pm}\exp[i(kx - \omega t)]
  \]
  for $A_{\pm} \in C$ and $k, \omega \in \R$.
  Such solutions obey the \textit{dispersion relation}:
  \[
    \omega = ck
  \]
  We call $\omega$ the \textit{angular frequency}.
  This is related to the \textit{wavelength} $\lambda$ and \textit{frequency} $\nu$ by:
  \[
    \nu = \frac{\omega}{2\pi} \text{ and } \lambda = \frac{2\pi c}{\omega} = \frac{c}{\nu}
  \]
  Waves can also \textit{interfere} with each other.
  If we have two waves with amplitudes $f_1$ and $f_2$, then $f_1 + f_2$ is a \textit{superposition} of the two waves.
  This interference can be \textit{constructive} (where the resulting wave is of larger amplitude), or \textit{destructive} (where the resulting wave is of a smaller amplitude).
\end{enumerate}
\section{Photoelectric Effect}
Suppose we have some monochromatic light (i.e. light of only one frequency) of frequency $\lambda$ incident on a metal plate.
When the photons interact with the electrons in the metal, they transfer energy to them, allowing them to escape from their atoms and be emitted from the plate.

We want to consider what happens as we increase the intensity $I = |A|^2$ (where $A$ is the \textit{amplitude}) of the light.

\subsection{Classical Expectation}
From a classical perspective, since incident light has $E \propto I$, we might expect to have to increase $I$ until there is enough energy to free the electrons from the metal.
We also expect the rate of emission of the electrons to be constant, even as we increase the intensity past the emission threshold.
Finally, we expect the velocity of emitted electrons to increase as $I$ (and therefore $E$) increases.
\subsection{Reality}
However, when we carry out this experiment we obtain the following three results:
\begin{enumerate}
  \item Below a certain frequency $\omega$, no electrons were emitted
  \item The velocity (and therefore the kinetic energy) of the electrons did not change with the intensity
  \item The emission rate increased with energy.
\end{enumerate}
In 1905, Einstein theorised an explanation.
He theorised that light comes in small \textit{quanta} called \textit{photons}, each carrying a small amount of energy:
\[
  E = \hbar \omega = \frac{h \omega}{2\pi}
\]
and momentum:
\[
  p = \hbar k
\]
where $h = \qty{6.63e-34}{\joule\second}$ is the \textit{Planck constant} and $\hbar = \frac{h}{2\pi} = \qty{1.055e-34}{\joule\second}$ is the \textit{reduced Planck constant}.

Electron emission is then explained by a single photon interacting with a single electron.
The kinetic energy of the emitted electron is then the difference between the binding energy $\phi$ of the electron and the energy of the incident photon:
\[
  E = \hbar \omega - \phi
\]
Thus, since $E \geq 0$, the minimum energy for an incident photon to cause an emission is $E_{\text{min}} = \phi$, or equivalently, $\omega_{\text{min}} = \phi / \hbar$.
This explains \textbf{(i)}, as we have to exceed the threshold frequency for emission, and also \textbf{(ii)}, as we see the energy of emitted electrons does not depend on the intensity.

Light quanta also explains explains \textbf{(iii)}.
Increasing intensity (or equivalently energy), leads to more incident light quanta per unit time (since they are monochromatic they all have the same energy), and therefore a greater emission rate of electrons provided $\omega \geq \omega_{\text{min}}$.
\section{Atomic Spectra}
In 1908, Rutherford supposed that the atom contained a small nucleus where the mass and positive charge of the atom was concentrated with electrons orbiting around it in a continuous cloud.
However this model was flawed as electrons moving in a circular orbit would radiate energy, but nothing was preventing them from collapsing in towards the nucleus.
Furthermore, when we irradiated an atom was light, it gained energy and then re-emitted light as it lost energy.
However the wavelength of this re-emitted light was only ever at set of discrete frequencies given by:
\[
  \omega_{mn} = 2\pi R_0 \left(\frac{1}{n^2} - \frac{1}{m^2}\right)
\]
where $n, m \in \N$ and $R_0$ is the \textit{Rydberg constant}.
This suggests that the electrons instead orbit at set energy levels.

In 1913, Bhor proposed that the orbits around the nucleus are quantised (i.e. electrons orbit at discrete levels) and that their angular momentum around the nucleus is given by:
\begin{equation}
  L_{n} = n \hbar \label{quantL}
\end{equation}
for each $n \in \N^{+}$.
\begin{proposition}
  If the angular momentum $L$ is quantised, then so are the radius $r$, velocity $v$, and energy $E$.
\end{proposition}
\begin{proof}
  The angular momentum of an electron at position $\vec{r}$ is given by:
  \[
    \vec{L} = m_e \vec{v} \times \vec{r}
  \]
  We assume that the orbit is circular with radius $r$ and so:
  \[
    L = |\vec{L}| = m_e vr \implies v = \frac{L}{m_e r}
  \]
  So if we have an electron in the $n$-th energy level, using \cref{quantL}, its velocity is:
  \[
    v_{n} = \frac{n \hbar}{m_e r}
  \]
  Thus the velocity is also quantised.

  We now want to use Newton's Second Law.
  The electron moves in a circular orbit under the Coulomb force for charge $q = e$, with centripetal acceleration $a = v^2 / r$ (see IA Dynamics and Relativity).
  Thus we have:
  \[
    F = - \frac{e^2}{4 \pi \varepsilon_0} \frac{1}{r^2} = m_e \frac{v^2}{r}
  \]
  Using our expression for $v_n$, we have:
  \[
    r_n = \frac{4\pi \varepsilon_0}{m_e e^2} {\hbar}^2 n^2
  \]
  and so the radius is quantised.
  Since $r_n$ increases with $n$, this gives us a minimum radius of:
  \[
    r_1 = \frac{4\pi \varepsilon_0}{m_e e^2} {\hbar}^2
  \]
  This is called the \textit{Bohr radius}.
  We can then write $r_n = r_1 n^2$.

  For the energy, we use the potential for the Columb force:
  \begin{align*}
    E &= \frac{1}{2} m_e v^2 - \frac{e^2}{4 \pi \varepsilon_0} \frac{1}{r} \\
    \implies E_n &= -\frac{e^2}{8 \pi \varepsilon_0 r_1} \frac{1}{n^2}
  \end{align*}
  and so energy is also quantised.
\end{proof}
So electrons exist at discrete \textit{energy levels} $E_n$.
They can interact with photons to gain energy, causing them to jump to a higher energy level, and then over time, decay back to a lower energy level, emitting a photon with energy equal to the difference between the energy levels.
We call the lowest energy level $E_1$, the \textit{ground state}.
\begin{remark}
  We can use the above to predict the Rydberg constant.
  We defined $\omega_{m n}$ as the energy of light emitted by photons as they go from energy level $m$ to $n$, thus using the above we have:
  \begin{align*}
    \omega_{mn} &= \frac{1}{\hbar}(E_m - E_n) \\
                &= 2\pi c \left(\frac{e^2}{4\pi \varepsilon_0 \hbar c}\right)^2 \left(\frac{1}{n^2} - \frac{1}{m^2}\right)
  \end{align*}
  and so:
  \[
    R_0 = \left(\frac{e^2}{4\pi \varepsilon_0 \hbar c}\right)^2
  \]
\end{remark}
\end{document}
